\section{Data}
\label{sec:data}

AmphionASR is initialised from the public
Qwen3-ASR-1.7B weights~\cite{qwen2025qwen3asr} and trained through a
three-stage LoRA supervised fine-tuning (SFT) recipe
over eight curated training sets
(Table~\ref{tab:training_sets}), totalling $3.74$\,million utterances
built from eleven public corpora (Table~\ref{tab:data_public}).
The sets map directly onto the recognition settings introduced in
Section~\ref{sec:intro}: whispered-speech ASR, degradation-robust
ASR, hotword-conditioned ASR, target-speaker ASR (TS-ASR), and
general transcription data included to reduce the risk of catastrophic
forgetting.
The eight sets are distributed across the three SFT stages:
degradation and whisper data in Stage~1, hotword and target-speaker
data in Stage~2, and a replay of the Stage~1--2 sets together with
general-transcription corpora in Stage~3. Section~\ref{sec:training}
documents the per-stage optimiser settings.

\begin{table}[t]
  \caption{Training sets used across the AmphionASR three-stage
    SFT recipe.
    ``Utterances'' counts training examples after offline preparation;
    TS-ASR rows contain two audio inputs (enrollment + mixture) per
    example.}
  \label{tab:training_sets}
  \centering
  \small
  \setlength{\tabcolsep}{3pt}
  \begin{tabular}{p{4.8cm}lrp{3.2cm}}
    \toprule
    Training set & Task & Utterances & Notes \\
    \midrule
    WhispEar \cite{fang2025whisperear}
      & Whispered-speech ASR & 222\,440  & bilingual whisper corpus \\
    Voice-in-the-Wild-2M \cite{xie2026mega}
      & Degradation     & 556\,421  & real-world degradations \\
    Five-corpus hotword mixture
      & Hotword ASR     & 1\,043\,317 & single-audio \\
    AISHELL-2 \cite{du2018aishell2}
      & General ASR      & 1\,009\,223 & Mandarin retention \\
    LibriSpeech \cite{panayotov2015librispeech}
      & General ASR      & 281\,241  & English retention \\
    CommonVoice~EN target-speaker
      & TS-ASR          & 184\,747  & synthesised \\
    MagicData target-speaker
      & TS-ASR          & 106\,116  & synthesised \\
    Multi-corpus target-speaker ASR
      & TS-ASR          & 333\,640  & synthesised \\
    \midrule
    \textbf{Total}            &                 & \textbf{3\,737\,145} & \\
    \bottomrule
  \end{tabular}
\end{table}

\begin{table}[t]
  \caption{Public corpora that appear in the training mixture.
    A corpus may contribute to multiple training sets (e.g.\
    CommonVoice~EN feeds both hotword-conditioned and target-speaker
    training); utterance counts are \emph{not} additive across rows.}
  \label{tab:data_public}
  \centering
  \small
  \setlength{\tabcolsep}{4pt}
  \begin{tabular}{p{4.5cm}p{6.8cm}}
    \toprule
    Corpus & Role in SFT \\
    \midrule
    WhispEar \cite{fang2025whisperear}
      & Whispered-speech ASR ($222$\,k utterances) \\
    Voice-in-the-Wild-2M \cite{xie2026mega}
      & Degradation-robust ASR ($556$\,k) \\
    CommonVoice EN (cleaned) \cite{ardila2020commonvoice}
      & Hotword-conditioned ASR ($500$\,k cap) + target-speaker syntheses \\
    MagicData Mandarin \cite{magicdata}
      & Hotword-conditioned ASR ($350$\,k cap) + target-speaker syntheses \\
    AISHELL-1 \cite{bu2017aishell}
      & Hotword-conditioned ASR \\
    AISHELL-2 \cite{du2018aishell2}
      & General Mandarin ASR retention ($1.01$\,M) \\
    AISHELL-3 \cite{shi2021aishell3}
      & Hotword-conditioned ASR + target-speaker synthesis \\
    LibriSpeech \cite{panayotov2015librispeech}
      & General English ASR retention (train-clean-100/360 + train-other-500) \\
    MLS English \cite{pratap2020mls}
      & Target-speaker synthesis \\
    THCHS-30 \cite{wang2015thchs30}
      & Hotword-conditioned ASR + target-speaker synthesis \\
    aidatatang\_200zh
      & Target-speaker synthesis \\
    \bottomrule
  \end{tabular}
\end{table}

\subsection{Whispered-Speech ASR}
\label{subsec:whisper-data}

The WhispEar training set covers the full bilingual training
pool ($222{,}440$ utterances): $199{,}805$ English
whisper clips from the Emilia-based subset, $18{,}588$ English
whisper clips from the WTIMIT subset, and $4{,}047$ Mandarin whisper
clips from the wEar subset~\cite{fang2025whisperear}.
Section~\ref{sec:whisper-asr} reports evaluation on the wEar and
WTIMIT whisper test sets.

\subsection{Degradation-Robust ASR}
\label{subsec:vitw-data}

The Voice-in-the-Wild-2M training set provides $556{,}421$
single-utterance examples covering far-field capture, additive noise,
reverberation, clipping, dropout and their superpositions~\cite{xie2026mega}.

\subsection{Hotword-Conditioned ASR}
\label{subsec:hotword-sources}

Hotword training data are materialised offline from five public
ASR corpora whose transcripts carry mined entity lists
(Table~\ref{tab:data_hw_train}).
At conversion time each utterance receives a prompt-level hotword line
with probability $0.8$; the line combines real hotwords with
hard-negative and random distractors, with the total list size
controlled in $[10, 50]$ (Section~\ref{subsec:hotword-pool}).
The resulting training set contains only single-audio examples.

\begin{table}[t]
  \caption{Public corpora in the hotword-conditioned training set.
    ``Cap'' is the utterance cap applied during offline preparation.}
  \label{tab:data_hw_train}
  \centering
  \small
  \begin{tabular}{lrr}
    \toprule
    Source & Cap & Utterances \\
    \midrule
    CommonVoice EN (cleaned) \cite{ardila2020commonvoice} & 500\,k & 500\,000 \\
    MagicData Mandarin \cite{magicdata}                     & 350\,k & 349\,958 \\
    AISHELL-1 \cite{bu2017aishell}                            & --     & 120\,098 \\
    AISHELL-3 \cite{shi2021aishell3}                          & --     & 63\,261 \\
    THCHS-30 \cite{wang2015thchs30}                            & --     & 10\,000 \\
    \midrule
    \textbf{Total}                                            &        & 1\,043\,317 \\
    \bottomrule
  \end{tabular}
\end{table}

Two CommonVoice hotword \emph{test} splits (ZH and EN) with the same
annotation protocol are reserved for evaluation
(Table~\ref{tab:data_hotwords_eval}).

\begin{table}[t]
  \caption{Hotword evaluation splits (not used in SFT).}
  \label{tab:data_hotwords_eval}
  \centering
  \small
  \begin{tabular}{ll}
    \toprule
    Split & Underlying audio \\
    \midrule
    CommonVoice ZH hotwords & CommonVoice ZH test \\
    CommonVoice EN hotwords & CommonVoice EN test \\
    \bottomrule
  \end{tabular}
\end{table}

\subsection{General Transcription Retention}
\label{subsec:general-asr-retention}

To limit catastrophic forgetting of general transcription behaviour, the
training mixture retains exposure to two large single-speaker corpora:
the full AISHELL-2 training split ($1{,}009{,}223$ utterances) and
LibriSpeech train-clean-100/360 plus train-other-500 ($281{,}241$
utterances).

\subsection{TS-ASR Synthesis}
\label{subsec:tsasr-synth}

Target-speaker training examples are produced by an offline
synthesiser. From single-speaker source manifests, each output sample
contains a target waveform mixed with $K$ interferers, an enrollment
utterance from the same speaker but \emph{a different recording}, and
the target's transcript; a fraction of \emph{anti-hallucination
negatives}---empty-transcript samples whose target speaker is removed
from the mixture---are intended to expose the model to cases where the
requested speaker is absent. The mixing equation, RIR/noise pools, and negative
sub-types are detailed in Appendix~\ref{sec:appendix-tsasr-synth}.

\paragraph{Target-speaker training sets.}
Stage~2 of the SFT recipe includes three synthesised
target-speaker partitions totalling $624{,}503$ examples
(Table~\ref{tab:ts_training_sets}).
The first two are built from single-corpus manifests (CommonVoice~EN
and MagicData respectively); the third draws from six public corpora.
Positive counts by underlying corpus in the multi-corpus partition
are: CommonVoice~EN ($225{,}638$), MLS English ($40{,}184$), MagicData
($10{,}119$), aidatatang\_200zh ($4{,}140$), AISHELL-3 ($1{,}425$),
and THCHS-30 ($483$); negatives account for the remaining
$51{,}651$ examples ($\approx\!15$\,\%).

\begin{table}[t]
  \caption{Target-speaker training sets in the training mixture.
    ``Pos.'' counts positive (target-speaker-present) examples;
    ``Neg.'' counts anti-hallucination negatives with empty
    transcripts.}
  \label{tab:ts_training_sets}
  \centering
  \small
  \setlength{\tabcolsep}{4pt}
  \begin{tabular}{llrr}
    \toprule
    Set & Primary source & Pos. & Neg. \\
    \midrule
    CommonVoice~EN & CommonVoice EN   & 175\,510 & 9\,237  \\
    MagicData      & MagicData        & 100\,810 & 5\,306  \\
    Multi-corpus   & 6 corpora        & 281\,989 & 51\,651 \\
    \midrule
    \textbf{Total} &                  & 558\,309 & 66\,194 \\
    \bottomrule
  \end{tabular}
\end{table}

Section~\ref{sec:tsasr} reports evaluation on an in-house test set
generated by the same pipeline with a held-out random seed.

\subsection{Hotword Pool Construction at Training Time}
\label{subsec:hotword-pool}

For hotword-bearing training sets, training-time hotword lists are
assembled per sample so that the total number of real hotwords and
distractors is controlled in $[10, 50]$. Each list is built in three
stages:

\begin{enumerate}
  \item \textbf{Real terms.} Each ground-truth hotword is
        independently dropped with probability $0.05$ to simulate
        imperfect retrieval recall. Keeping most real hotwords lets
        the model use them as contextual bias to correct rare entities
        and domain-specific terms that an unconditioned decoder would
        mis-transcribe.
  \item \textbf{Hard negatives.} An online retriever fetches up to
        $10$ confusable terms per sample using the weighted
        score $0.6\cdot\text{char-overlap} +
        0.4\cdot\text{bigram-Jaccard}$ over a per-corpus hotword
        index, augmented for Mandarin by a syllable-level
        nearest-neighbour search. These phonetically or
        orthographically similar distractors are intended to expose the
        model to distinctions among look-alike hotwords rather than
        favouring any listed term.
  \item \textbf{Random distractors.} The remaining slots are filled
        from the global hotword pool, with the hard-to-random ratio
        set so that hard negatives saturate first and the list total
        stays within $[10, 50]$. Random fillers expose the model to
        irrelevant hotwords and control the prompt list length; this is
        intended to discourage reliance on non-matching candidates.
\end{enumerate}

The whole hotword prompt line is then kept with probability $0.8$
and dropped otherwise. This exposes the model to inputs both with and
without a hotword side channel during training.
